\chapter{Caching and Reuse}
\label{ch:caching}

\begin{chapterintro}
Caching keeps a result Spark would otherwise recompute --- but the wrong cache
wastes storage memory and starves execution, and a cache used once is pure
overhead. Spark also reuses work you did not ask it to: an identical exchange or
subquery computed once and read twice. This chapter covers when caching pays,
which storage level to pick against the GC budget from \Cref{ch:execmem}, and how
to see reuse in a plan. It is also where execution and storage memory finally
meet at runtime.
\end{chapterintro}

\section{When caching pays}
\tbd

\section{Lazy \texttt{cache()} versus eager \texttt{CACHE TABLE}}
\tbd

\section{Storage levels and the GC trade-off}
\tbd

\section{How execution and storage memory borrow at runtime}
\tbd

\section{Exchange and subquery reuse}
\tbd

\section{Unpersist and cache thrashing}
\tbd

\section{Summary}
\tbd

\exercisehead
\begin{exercises}
\item Build a query that reads one CTE twice. Show the plan reusing the exchange,
      then change the query so the reuse breaks and measure the difference.
\end{exercises}
